% Generated draft: requires mathematical and source review.
\documentclass[11pt]{article}
\usepackage[margin=25mm]{geometry}
\usepackage{fontspec}
\setmainfont{Latin Modern Roman}
\usepackage{amsmath,amssymb,mathtools,unicode-math}
\setmathfont{Latin Modern Math}
\usepackage{enumitem,needspace,xurl,hyperref}
\setlength{\parindent}{0pt}
\setlength{\parskip}{5pt}
\setlength{\emergencystretch}{3em}
\newcommand{\N}{\mathbb{N}}
\newcommand{\Z}{\mathbb{Z}}
\newcommand{\Q}{\mathbb{Q}}
\newcommand{\R}{\mathbb{R}}
\newcommand{\C}{\mathbb{C}}
\newcommand{\F}{\mathbb{F}}
\newcommand{\A}{\mathbb{A}}
\newcommand{\PP}{\mathbb{P}}
\newcommand{\E}{\mathbb{E}}
\newcommand{\eps}{\varepsilon}
\newcommand{\dd}{\,\mathrm{d}}
\begin{document}
\section*{UnsolvedMath ID: 30005673}
\textbf{Source problem code:} catalogue entry
\subsection{Definitions and mathematical statement}
Let \(G\) belong to the class of groups intended by the source, and let \(M\) be a compact smooth \(G\)-manifold. The supplied record does not specify this class of groups, whether \(M\) may have boundary, or the boundary conventions. These hypotheses require clarification before the statement has a fully specified domain.

Write \(M_+=M\sqcup\{*\}\), with a disjoint \(G\)-fixed basepoint, and let \(\Sigma_G^\infty M_+\) denote its equivariant suspension spectrum. The notation \(A_G(M)\) denotes the equivariant algebraic \(K\)-theory of spaces referred to in the source, whose construction uses spectral Mackey functors. The intended category of equivariant spectra and its indexing universe are not specified in the supplied excerpt.

The proposed equivariant splitting asks whether, for every such \(M\), there is an equivariant Whitehead spectrum \(\operatorname{Wh}_G(M)\) and an equivalence of equivariant spectra
\[
 A_G(M)\simeq \Sigma_G^\infty M_+\times\operatorname{Wh}_G(M),
\]
with an equivalence of spaces
\[
 \Omega^{\infty+1}\bigl(\operatorname{Wh}_G(M)^G\bigr)
 \simeq H_G^\infty(M).
\]
Here the superscript \(G\) denotes the intended fixed-point spectrum, and \(\Omega^{\infty+1}E\) means \(\Omega\Omega^\infty E\). The fixed-point construction requires specification. The space \(H_G^\infty(M)\) is intended to be the geometrically defined stable equivariant \(h\)-cobordism space of \(M\), stabilized using representation disks; it is not being defined by the displayed equivalence. Its precise construction, stabilization maps, and indexing representations are absent from the record. The source also does not specify naturality or compatibility requirements for the splitting.

An equivariant \(h\)-cobordism \((W;M,N)\) is an \(h\)-cobordism between compact \(G\)-manifolds for which the boundary inclusions \(M\hookrightarrow W\) and \(N\hookrightarrow W\) are \(G\)-homotopy equivalences, meaning that their homotopy inverses and homotopies are equivariant. For a finite-dimensional real \(G\)-representation \(V\) equipped with an invariant inner product, write \(D(V)\) for its unit disk. The source emphasizes stabilization by products with such disks, but does not supply all hypotheses ensuring this representation convention or the necessary boundary and corner conventions.

As background motivating this stabilization, the supplied excerpt attributes to its reference [13] the assertion that the equivariant Whitehead torsion of \(M\hookrightarrow W\) vanishes if and only if there is a representation \(V\) for which
\[
 (W\times D(V);\,M\times D(V),\,N\times D(V))
\]
is a trivial equivariant \(h\)-cobordism. This is a source-reported assertion with unspecified hypotheses and an unidentified reference [13], not an additional theorem asserted here. The equivariant Whitehead group occurring in that assertion must be distinguished from the spectrum \(\operatorname{Wh}_G(M)\) in the conjecture.

The excerpt further discusses possible dimension differences between fixed-point strata and isovariant variants. An equivariant map \(f:X\to Y\) is isovariant when \(G_x=G_{f(x)}\) for every \(x\in X\), where \(G_x=\{g\in G:gx=x\}\). An equivariant \(h\)-cobordism is called isovariant there when its boundary inclusions are isovariant homotopy equivalences. The excerpt terminates before completing its subsequent hypotheses; consequently no particular dimension-gap condition or restriction to isovariant cobordisms can be imposed from the supplied evidence.
\subsection{Short English statement}
The question asks whether the equivariant \(A\)-theory spectrum of every compact smooth \(G\)-manifold splits into its equivariant suspension spectrum and a Whitehead spectrum whose fixed points recover stable equivariant \(h\)-cobordisms after applying \(\Omega^{\infty+1}\). Stabilization must use representation disks. The group class, boundary conventions, spectrum and fixed-point models, and geometric stabilization construction remain unspecified. The supplied label is partially solved; this draft does not independently verify that assessment.
\subsection{Sources}
\begin{itemize}
\item Supplied local problem-registry record, permanent UnsolvedMath ID \(30005673\), exact source problem code \texttt{catalogue entry}, title \emph{Equivariant Splitting of A-Theory}. The record gives proposed year 2023 and no named proposer.
\item Supplied report citation: \emph{Homotopy Theory (2024)}, Oberwolfach Reports, pp.\ 1893--1952, DOI \texttt{10.4171/owr/2023/34}. Supplied source URL: \url{https://doi.org/10.4171/owr/2023/34}. The citation's 2024 label and the record's 2023 year are retained without reconciliation.
\item Malkiewich--Merling, \emph{Equivariant parametrized h-cobordism}, supplied literature reference: \url{https://doi.org/10.1016/j.aim.2022.108242}.
\item Goodwillie--Igusa--Malkiewich--Merling, \emph{h-cobordism functoriality}, supplied literature reference: \url{https://arxiv.org/abs/2303.14892}.
\item \emph{OWR Homotopy Theory report}, supplied literature reference: \url{https://doi.org/10.4171/OWR/2023/34}.
\item Literature assessment dated exactly \texttt{2026-08-22}: the supplied curation labels the problem partially solved and reports a map from equivariant \(A\)-theory to the suspension spectrum with fiber described by stable \(h\)-cobordism spectra, known stable equivariant \(h\)-cobordism splittings, and an open full \(A\)-theory product splitting for every compact smooth \(G\)-manifold. This is supplied dated curation derived from the OWR material, not an independently verified literature assessment; its claims and time-sensitive status require independent checking.
\item The registry labels the statement \texttt{corrected\_verified} and reports that it was reconstructed and checked against the official Oberwolfach report. That verification is supplied provenance, not verification performed for this draft. The excerpt cites [14] for the equivariant \(A\)-theory construction, [6, 7, 10, 5, 4] for spectral Mackey functors, and [13] for the torsion assertion, but supplies no bibliography identifying these numbered references.
\end{itemize}
\end{document}
